\section{Introduction}
\label{sec:introduction}

The \texttt{wsq-viewer} project needs a WSQ codec that runs inside Visual Studio Code, without
Python or native dependencies, fast enough that opening a fingerprint feels instantaneous, and
correct enough to be trusted for forensic images. The companion document \emph{The WSQ
Fingerprint Image Compression Algorithm} describes the algorithm; this one describes how its
implementation was made fast and how fast it is.

The codec is a Rust library (\code{crates/wsq}) with a command-line tool and a WebAssembly
build for the extension. It has two arithmetic contracts:
\begin{description}[leftmargin=0pt]
  \item[Exact mode] reproduces NBIS 5.0.0, the NIST reference implementation, bit for bit: the
    decoder returns exactly the pixels of \code{dwsq}, and the encoder writes exactly the bytes of
    \code{cwsq}, comment included.
  \item[Fast mode] lets the arithmetic run in any order, with fused multiply-add and double
    precision where cheaper, and stays within the compliance tolerances of the specification.
\end{description}
Both modes run on any of five instruction sets, selected at run time.

Table~\ref{tab:summary} summarizes the outcome on a 1452$\times$1006 four-finger slap
compressed at 0.75 bits per pixel, on one performance core of a 2022 laptop processor. The rest
of the report explains how these numbers were obtained (Sections~\ref{sec:method}
to~\ref{sec:optimizations}), presents the measurements on the whole NIST data set
(Section~\ref{sec:results}), answers whether giving up exactness would buy more speed
(Section~\ref{sec:relaxation}), and reports a test of batch decoding on a GPU
(Section~\ref{sec:gpu}).

% Generated by tools/gen_tables.py; do not edit by hand.
\begin{table}[ht]
  \centering\small
  \begin{tabular}{@{}lrrrr@{}}
    \toprule
    & \multicolumn{2}{c}{decode} & \multicolumn{2}{c}{encode}\\
    \cmidrule(lr){2-3}\cmidrule(l){4-5}
    implementation & ms & vs NBIS & ms & vs NBIS\\
    \midrule
    NBIS 5.0.0 (C), one core & 52.38 & 1.0$\times$ & 36.74 & 1.0$\times$\\
    scalar, exact & 9.77 & 5.4$\times$ & -- & --\\
    SSE2, exact & 7.15 & 7.3$\times$ & -- & --\\
    AVX2, exact & 4.80 & 10.9$\times$ & 6.94 & 5.3$\times$\\
    AVX2, fast & 4.49 & 11.7$\times$ & 7.09 & 5.2$\times$\\
    AVX2, exact, 4 threads & 2.66 & 19.7$\times$ & -- & --\\
    WebAssembly in Node.js, scalar, exact & 18.67 & 2.8$\times$ & -- & --\\
    WebAssembly in Node.js, SIMD, exact & 6.50 & 8.1$\times$ & -- & --\\
    \bottomrule
  \end{tabular}
  \caption{Image \texttt{a002} of the NIST set (1452$\times$1006, a four-finger slap) at 0.75 bits per pixel: median times. Exact means the same pixels and bytes as NBIS.}
  \label{tab:summary}
\end{table}

